% 本轮核查：与译文并行保存，不将编译成功等同于数学审定。
\head{第七版：19.5--19.6续译及校订正文}
\section*{范围与采用的编辑方式}
\addcontentsline{toc}{section}{范围与采用的编辑方式}
本轮以第六版实际工程为起点，从第0章（19.5.1）连续译至（19.6.7）末尾。底本为上传的\texttt{EGAIV.1.pdf}，上传PDF第89--103页，分册印刷第90--104页。其中第99页跨19.5和19.6，第103页在（19.6.7）后进入19.7；后者没有算入已译范围。新增内容包括19.5的一般定理及其全部分段证明，以及19.6的Cohen定理、系数域、正则性与形式光滑性判据。

本版执行“正文写校正后的内容，原文错误在脚注注明”。因此，正文中的必要条件、已核实的源靶及可靠补证直接修正；脚注说明原来的写法、为什么修改，以及属于作者勘误还是译校补证。已经撤销、而本轮没有取得可靠替代定理的条目，只保留编号与撤销状态。旧稿的回溯修改限于本轮直接复核的（19.3.5，v）、（19.3.10.1）、（19.3.12）及（19.4.1）；\textbf{不声称已把此前约九百页内全部历史校订意见逐项改写完成}。

新增注号以G7开头。以下区分原命题的反例、证明步骤的反例、文字错误和使用范围提醒。同一条证明中的多个排印问题不计作多个错误定理；本次补证是译校分析，不冒称作者原证。

\begin{longtable}{@{}p{.16\linewidth}p{.21\linewidth}p{.54\linewidth}@{}}
\toprule 编号 & 分类 & 核查结论\\\midrule\endhead
G7-E01 & 条件遗漏，有反例 & （19.6.2）须给定相容的 $k\subset A$，只给 $k\subset K$ 不够。\\
G7-G01 & 证明步骤错误，可补 & （19.5.6.7）不能比较原同态核与相伴分次核；以有限滤过逐级消去补证。\\
G7-G02 & 证明步骤错误，可补 & （19.6.6）扩域后的分量不必同构于原环；改用Galois作用的传递性与下降。\\
G7-H01 & 拓扑条件需明示 & （19.5.1）的商模不在任意拓扑情形下自动离散；后文应用具备所需条件。\\
G7-T01--T08 & 对象、次数、引文等 & 分次生成对象、指标、公式号、正次数上同调、零元、张量积对象、假设依赖。\\
G7-H02--H05 & 使用范围与核对 & 完备化、增广商环、相容提升的理想像条件、Zariski环约定。\\
\bottomrule
\end{longtable}

\audit{G7-E01}{（19.6.2）：剩余域中的子域不自动给出环中的子域}{原条件按字面不足，正文已补明嵌入}
\oldauditsrc{EGA IV第一部分，印刷第101页，上传PDF第100页}
原文设 $A$ 为分离完备局部环、$K$ 为其剩余域，只给出 $k\subset K$ 且 $K/k$ 可分，却断言存在包含 $k$ 的子域 $K'\subset A$，使剩余同态在 $K'$ 上同构到 $K$。原证明随即把 $A$ 当成 $k$-代数，并以 $k$-同态作提升；但所列假设没有提供这一结构。

\lead{反例}
固定素数 $p$，取
\[
 A=\mathbf Z_p=\varprojlim_n\mathbf Z/p^n\mathbf Z,
 \qquad K=\mathbf F_p,\qquad k=\mathbf F_p.
\]
$A$ 对其极大理想 $(p)$ 的进制拓扑分离完备，剩余域为 $K$；恒等扩张 $K/k$ 可分。因此初版所有明确列出的条件均满足。
若存在保单位的域嵌入 $i:\mathbf F_p\to\mathbf Z_p$，则
\[
 0=i(p\cdot1_{\mathbf F_p})=p\cdot1_{\mathbf Z_p},
\]
但 $p\cdot1_{\mathbf Z_p}$ 的模 $p^2$ 分量为非零的 $p\bmod p^2$，矛盾。因此该环甚至不含同构于 $k$ 的子域，更不可能有原文所述的 $K'$。

\lead{修订版本及证明}
必须给定相容的保单位嵌入 $k\hookrightarrow A$，并用剩余映射把 $k$ 识别为 $K$ 的子域。由 $K/k$ 可分及（19.6.1），$K$ 为形式光滑的离散 $k$-代数。对 $A$ 及 $\mathfrak m$ 用（19.3.10），可把 $K\xrightarrow{\mathrm{id}}A/\mathfrak m$ 连续提升为 $k$-同态 $s:K\to A$。它因 $K$ 为域而单射，且 $\pi s=\mathrm{id}_K$。所以 $K'=s(K)$ 就是所求子域；$s$ 为 $k$-同态保证 $k\subset K'$。

这项补充并不把混合特征环错误地变成等特征环，也没有断言系数域对所有完备局部环都存在。紧接的（19.6.3）使用的是已嵌入 $A$ 的素子域，故该处满足修订条件。

\audit{G7-G01}{（19.5.6.7）：从相伴分次信息回到原核须逐级处理}{证明中的一般包含不成立；本版给出替代证明}
\oldauditsrc{EGA IV第一部分，印刷第98页，上传PDF第97页}
初版在证明后半部使用类似
\[
 \Ker h\subseteq\Ker(\gr h)
\]
的包含，将已分次的源识别为自身的相伴分次环。对于保持递降滤过的同态，这并不是一般性质：某个元素的最低齐次项可以映到更高滤过层，并与该元素的高次部分抵消。只能断言其\emph{首项}被相应的分次同态杀掉，不能断言整个元素被 $\gr h$ 杀掉。

\lead{具体检验例}
设 $k$ 为任意域，令
\[
 R=k[x,y]/(x,y)^3,\qquad S=k[t]/(t^3),
\]
分别取极大理想幂滤过。定义保滤过 $k$-代数同态
\[
 h:R\to S,\qquad h(x)=t,\quad h(y)=t^2.
\]
它良定义，因为任何总次数至少3的单项式均映到 $(t^3)=0$。源、目标本来就是分次环，所以可与其相伴分次环识别。取
\[
 z=y-x^2.
\]
则 $h(z)=t^2-t^2=0$，但 $\gr h$ 的一次分量把 $y$ 送到零、把 $x$ 送到 $t$，二次分量把 $x^2$ 送到 $t^2$，故
\[
 (\gr h)(z)=-t^2\ne0.
\]
因此所用包含按一般性质解释是错误的。本例\textbf{只反驳该证明步骤}，不是满足（19.5.6.7）全部特设条件而反驳其结论的例子。

\lead{正确的有限滤过消核引理}
固定截断次数 $n$。沿用正文的记号，令
\[
 E_{\alpha,n}=B/(\mathfrak b_\alpha+\mathfrak j^{n+1}),\qquad
 F_{\alpha,n}=\Salg_{C_\alpha}
 (\mathfrak j_\alpha/\mathfrak j_\alpha^2)/(F^+)^{n+1}.
\]
$F_{\alpha,n}$ 有次数 $0,\ldots,n$ 的分次，$F_{\alpha,n}^{\geq r}$ 是其增广理想的 $r$ 次幂。记 $E$ 系的过渡同态为 $p_{\alpha\beta}$，$F$ 系的为 $q_{\alpha\beta}$。给定正文中的满同态
$v_0:F_{\beta_0,n}\twoheadrightarrow E_{\alpha_0,n}$，其中 $\alpha_0=\alpha$、$\beta_0=\beta$。

由每个 $\varphi_r$ 为形式双态射和（19.1.3），对任意已选 $\beta_i$，可取一个共同的 $\alpha_{i+1}\geq\beta_i$，使对所有 $0\leq r\leq n$，过渡同态都有分解
\[
 q_{\beta_i\alpha_{i+1}}^{\,r}
 =w_i^r\varphi_{\alpha_{i+1}\alpha_{i+1},r}.
\]
共同指标的存在只用到 $n+1$ 个次数的有限性。再由（19.5.6.3，ii）取
$\beta_{i+1}\geq\alpha_{i+1}$ 和满同态
$v_{i+1}:F_{\beta_{i+1},n}\to E_{\alpha_{i+1},n}$，使
\[
 p_{\alpha_i\alpha_{i+1}}v_{i+1}
   =v_iq_{\beta_i\beta_{i+1}}.
\]
这些选取重复 $N=n+1$ 次即可，\emph{没有要求选择无限相容系统}。

若 $x\in\Ker v_{i+1}\cap F_{\beta_{i+1},n}^{\geq r}$，记其次数 $r$ 的类为 $x_r$。由 $v_{i+1}(x)=0$ 可知
\[
 \varphi_{\alpha_{i+1}\alpha_{i+1},r}
       q_{\alpha_{i+1}\beta_{i+1}}^{\,r}(x_r)=0.
\]
应用 $w_i^r$ 便得
$q_{\beta_i\beta_{i+1}}^{\,r}(x_r)=0$，因而
$q_{\beta_i\beta_{i+1}}(x)\in F_{\beta_i,n}^{\geq r+1}$。
上述交换关系又给出
$v_iq_{\beta_i\beta_{i+1}}(x)=0$。因此每退回一步指标，核中元素的滤过次数至少提高1。

从 $\Ker v_N\subset F_{\beta_N,n}^{\geq0}$ 开始连续退回 $n+1$ 步，得到
\[
 q_{\beta_0\beta_N}(\Ker v_N)
       \subset F_{\beta_0,n}^{\geq n+1}=0.
\]
故 $q_{\beta_0\beta_N}$ 通过满同态 $v_N$ 唯一分解为一个 $A$-代数同态
$u:E_{\alpha_N,n}\to F_{\beta_0,n}$。交换关系及 $v_N$ 满射给出
\[
 v_0u=p_{\alpha_0\alpha_N}.
\]
取 $\gamma=\alpha_N$，对任意 $\mu\geq\gamma$ 使用 $u\circ p_{\gamma\mu}$，正好得到（19.5.6.7）的结论。

\lead{对19.5总定理的作用}
这一证明保留原命题的全部假设，未追加诺特性或完备性。它只替换原证明不成立的核比较步骤。因而后面的（19.5.6.8）仍可据此构造实际代数同态，再完成（19.5.3，iii）的证明。其正确性并非由小规模代数检验代替，而由上述任意 $n$ 的论证承担。

\audit{G7-G02}{（19.6.6）：Galois扩域后的局部分量彼此同构，不必同构于原环}{结论可保留；正文以传递作用和下降修补}
\oldauditsrc{EGA IV第一部分，印刷第104页，上传PDF第103页}
原证明在 $k'/k'_0$ 为有限Galois扩张、$A'_0=\widehat A\otimes_k k'_0$ 为完备局部环的情形，把 $A'=A'_0\otimes_{k'_0}k'$ 写成若干“与 $A'_0$ 同构”的局部环的直积。这个说法一般不成立。取
\[
 A=k=k'_0=\mathbf R,\qquad k'=\mathbf C.
\]
则 $A'_0=\mathbf R$，$A'=\mathbf C$ 只有一个局部分量，但它不与 $\mathbf R$ 为 $\mathbf R$-代数同构：作为向量空间，二者维数分别为2和1。这个例子仍满足 $A$ 的正则性和形式光滑性，因此只是原证明中的辅助断言错误。

\lead{应使用的分量传递性}
设 $l\subset L$ 为有限Galois扩张，$B_0$ 为完备诺特局部 $l$-代数，令
$D=B_0\otimes_l L$。它为有限 $B_0$-代数，因而完备半局部，并分解为有限个完备局部环的直积。令 $G=\operatorname{Gal}(L/l)$。$G$ 通过 $1\otimes\sigma$ 作用于 $D$，且
\[
 D^G=B_0.
\]
这一等式可直接用向量空间的正合性验证：$L^G$ 是映射
$L\to\bigoplus_{\sigma\in G}L$，$a\mapsto(\sigma(a)-a)_\sigma$ 的核；在 $l$ 上张量 $B_0$ 保持这个核，且 $L^G=l$。

$G$ 置换直因子的单位幂等元。若作用有两个以上轨道，取一个非空真轨道，其单位幂等元之和 $e$ 属于 $D^G=B_0$，且 $e\ne0,1$。但局部环无非平凡幂等元：$e(1-e)=0$，$e$ 与 $1-e$ 至少一个是单位，由此只能有 $e=0$ 或1。矛盾。因此作用在分量上是传递的。

传递性给出分量之间的 $B_0$-代数同构；相对于 $L$，这些同构可能是 $\sigma$-半线性的。半线性同构依然保持环的正则性，而且保持剩余域在 $L$ 上的可分性，因为它只将基域嵌入与一个域自同构复合。

\lead{从一个良好分量推出形式光滑性}
沿用命题（19.6.6）的 $b')$。将给定 $k'$ 放入有限准Galois扩张 $L/k$，令 $l/k$ 为其中的有限纯不可分部分，使 $L/l$ 为有限Galois扩张。良好分量 $K'[[T_1,\ldots,T_n]]$ 经有限扩域到 $L$ 后变为
\[
 \bigl(K'\otimes_{k'}L\bigr)[[T_1,\ldots,T_n]],
\]
这一步用 $L/k'$ 有限维，保证与系数的任意乘积可交换。由 $K'/k'$ 可分，系数代数为有限个在 $L$ 上可分的域的直积。故 $D=\widehat A\otimes_kL$ 至少有一个正则、且剩余域在 $L$ 上可分的局部分量。

$B_0=\widehat A\otimes_kl$ 为局部环，因为 $l/k$ 纯不可分。以上传递性使所有分量都有这些性质，故由（19.6.4）都在 $L$ 上形式光滑，其有限乘积 $D$ 也形式光滑（19.3.5，v）。

最后，$\widehat A$ 作为拓扑 $k$-向量空间形式投射：对离散向量空间的连续映射通过开子空间商，而每个向量空间都是投射模；从该商取线性提升再复合，提升依然连续。由于 $L/k$ 有限自由且忠实平坦，\textbf{应使用（19.4.6）的条件 $2^\circ$}，将 $D$ 的形式光滑性下降至 $\widehat A$，再由（19.3.6）得到 $A$ 形式光滑。初版此处援引（19.4.7），该条是模去定义理想的判据，不是这里的有限忠实平坦扩域判据。

以上补证没有援引尚在后文的（22.5.8），也没有通过循环使用（19.6.6）的待证结论来证明自身。

\audit{G7-H01}{（19.5.1）：代数商同构与离散商拓扑必须区分}{一般拓扑下需限定；后文用途不受影响}
\oldauditsrc{EGA IV第一部分，印刷第90页，上传PDF第89页}
代数公式
\[
 \Salg_C^n(V)/(U\Salg_C^{n-1}(V))\simeq\Salg_C^n(V/U)
 \qquad(n\geq1)
\]
始终成立。但按原文在对称代数上给出的拓扑，不能从 $U$ 在 $V$ 中开就自动断言左边离散。

取 $C=\mathbf Z_p$，给
\[
 V=\bigoplus_{r\geq1}\mathbf Z_p/p^r\mathbf Z_p
\]
配备离散拓扑，令 $U=0$。这是拓扑 $C$-模：每个固定向量仅有有限个非零分量，因而被某个开理想 $p^N C$ 零化；利用 $V$ 离散，数乘在每个点连续。

在对称代数定义的零点基本邻域中取 $U=0$，与一次部分 $V$ 相交得到 $p^N V$。对每个 $N$，$p^N V\ne0$，因为第 $N+1$ 个直和分量中的 $p^N$ 非零。因此诱导在 $\Salg_C^1(V)=V$ 上的拓扑不是离散拓扑，即使原给定的 $V$ 拓扑是离散的。

若存在开理想 $\mathfrak a$ 满足 $\mathfrak a V\subset U$，则在每个正次数上
$\mathfrak a\Salg_C^n(V)\subset U\Salg_C^{n-1}(V)$，该商就离散。原文前面讨论的“$V$ 的拓扑不比由 $C$ 诱导的模拓扑更细”保证这一点。正文把离散性明确限定到该情形；把括号“离散”读成一般断言时，上例说明其不成立。后文 $V=\mathfrak j/\mathfrak j^2$ 的应用正满足原文所说明的拓扑条件。

\section*{记号、引文与假设依赖的逐项记录}
\addcontentsline{toc}{section}{记号、引文与假设依赖的逐项记录}
\audit{G7-T01}{（19.5.4.2）：检验复合同态须用源代数的生成性}{对象写法修正}
复合从 $\gr^\bullet(F_n)$ 经 $\gr^\bullet(E_n)$ 返回 $\gr^\bullet(F_n)$，因而应根据\emph{源}由零次和一次部分生成，判定复合在每次为恒等。原文在此说 $\gr(E_n)$ 由一次部分生成，虽然该性质本身也对，却不是所写论证需要的对象。正文写明正确的源，并保留原差异于脚注。

\audit{G7-T02}{（19.5.6.5）：目标增广理想和同态下标}{类型校正}
$w'_\lambda$ 按定义取值于 $F_{\beta,n}$，所以其将 $\mathfrak j_\lambda$ 映入的是 $F_{\beta,n}^+$，不是原文该处所写的 $F_{\delta,n}^+$。末尾一次同态的下标相应统一为 $w_{\beta\lambda}$。这类下标修正不改变构造。

\audit{G7-T03}{（19.5.6.6）：相同的分次生成对象问题}{类型校正}
与G7-T01相同，待识别的复合以 $\gr^\bullet(F_{\lambda,n})$ 为源。正文利用它的零次和一次生成性，而不是把中间的 $\gr^\bullet(E_{\lambda,n})$ 误写为所需源。

\audit{G7-T04}{（19.5.7，iii）：公式号误引}{已对照原页更正}
原文将当前同态引作（19.5.4.1），应为（19.5.7.1）。同时，对照页面可确认该段印的是“投射 $C_\lambda$-模”；自动提取文本丢掉下标时，不应据此另报“原文把 $C_\lambda$ 写成 $C$”这一错误。

\audit{G7-T05}{（19.6.1）：Hochschild上同调消失仅指正次数}{次数范围明示}
当 $K/k$ 有限可分时，$K$ 是 $K\otimes_kK$ 的投射模，故 $\Ext^i(K,L)=0$ 对 $i>0$ 成立；$i=0$ 则是 $\Hom(K,L)$，一般不为零。最简例 $K=L=k$ 给出 $H_k^0(k,k)=k$。正文限定 $i>0$；扩张分类只需要 $i=2$。

\audit{G7-T06}{（19.6.3）：可逆性须排除零元}{边界用语修正}
原证明说 $A'$ 中每个元素可逆；实际应为每个非零元素。其原因是 $A'\cap\mathfrak m=0$，而不是两个含零子集的交为空。证明结构和结论均不变。

\audit{G7-T07}{（19.6.5.1）：可分性要求张量积约化}{结论对象校正}
从 $A\otimes_kk'$ 正则约化得到的是子环 $K\otimes_kk'$ 约化。原文末句简写成 $K$ 约化不足以证明可分性，因为任意域都约化，包括不可分扩张中的大域。正文保留完整对象。

\audit{G7-T08}{（19.6.7，i）：有限根式重数假设的用途写反}{依证明实际依赖修正}
原文说有限根式重数仅用于 $b')\Rightarrow a)$。修补后的该方向不需要它；反而 $c)\Rightarrow b)$ 明确使用假设所给的 $k_0$，以保证有限根式扩域后剩余域可分。正文据此改写备注，并保留下一条“以后证明 $a)$ 与 $c)$ 无需该假设”的历史说明。本轮未把尚未译入的（22.5.8）算作已证明或已翻译。

\audit{G7-H02}{（19.5.4.4）和（19.6.4）：原环与完备化}{核对正确，不另判原版错误}
（19.5.4.4）中由形式光滑性提升系数环，目标为 $\widehat B$；扫描原页确有帽号。类似地，（19.6.4）的形式幂级数结构说的是 $\widehat A$。本版保留帽号，不因提取文本丢失它而把原文判错。

\audit{G7-H03}{（19.5.5）：增广商环隐含的理想包含}{约定展开}
写 $E/\mathfrak b$ 为经常数项增广到 $C$ 的商代数，须有 $\mathfrak b\subset\mathfrak n$；再结合 $\mathfrak n^k\subset\mathfrak b$。正文把这一原有增广条件的含义明写出来，不增加无关的限制。

\audit{G7-H04}{（19.5.6.3）仍可应用已修订（19.3.10.1）}{应用条件逐项复核}
在此过渡同态 $E_{\gamma,n}\to E_{\alpha,n}$ 下，理想
$\mathfrak j_\gamma/\mathfrak j_\gamma^{n+1}$ 的像恰为
$\mathfrak j_\alpha/\mathfrak j_\alpha^{n+1}$，因为它们都由同一 $\mathfrak j\subset B$ 取像得到。两个理想幂零，故（19.3.10.1）所补入的理想像相等条件满足。这样才可确保提升同时保留指定的剩余同态；不能仅有普通的无约束提升。

\audit{G7-H05}{（19.5.8）：Zariski环与前置拓扑假设}{使用范围提醒，不列为新反例}
该注应连同（19.5.7）的假设及原著采用的Zariski环约定使用。有限表现性和完成函子的判据有其诺特背景，不能把“Zariski环”擅自换成任意线性拓扑环。正文未删除这些条件。

\section*{回溯修改与保留的历史记录}
\addcontentsline{toc}{section}{回溯修改与保留的历史记录}
（19.3.5，v）现已采用正交幂等元提升证明；（19.3.10.1）在正文补入 $f(\mathfrak k)=\mathfrak k'$ 并明确同时提升两条映射；（19.4.1）在正文同时选取使 $\mathfrak j_\alpha C=0$ 的指标；（19.3.12）正文不再出现一整段被撤销的命题和错误证明，其原陈述在脚注注明，完整历史译文保存于下页附录。相应旧注F6编号仍保留，便于追踪。

旧核查报告作为历史档案继续收入。对于本轮已修改的条目，旧报告中“正文仍保留原错误写法”的描述只代表第六版状态。新正文、G7分析及实际修改记录优先。工程中附\texttt{history/V7\_body\_corrections.diff}，可逐行比较这两份既有正文源文件。

\section*{核查边界与技术验证的含义}
\addcontentsline{toc}{section}{核查边界与技术验证的含义}
本轮对新增范围进行了原页对照、公式与证明类型检查，并针对上述问题给出反例或补证；不是对全部历史译文的再次逐句审定，也不是机械形式化验证。新译文及校订仍属工作稿，尤其本轮补证属于可供独立复核的数学论证。

编译、目录及交叉引用检查、实际页数、原条目保留检查和选定页面的渲染检查，均记录在工程的\texttt{checks/validation.json}及\texttt{checks/REVIEW.md}。这些检查回答文件是否可编译和版面是否可读，不代替数学正确性审定。报告不预先填写尚未实测的页数或比较结果。

\section*{附录：已撤销（19.3.12）的历史译文}
\addcontentsline{toc}{section}{附录：已撤销（19.3.12）的历史译文}
以下仅保存第六版中的原文对照记录，全部陈述和推理都属于被撤销的条目，不是校订正文中的有效命题。本版正文只保留编号、撤销状态及脚注。
\subsection*{历史条目（19.3.12）}
\noindent\textbf{作者勘误提示。}作者在勘误清单三（Err$_{\mathrm{IV}}$，8）明确要求删除本条。以下保留初版陈述及证明的中文译文，仅供原文对照，不将其作为有效定理引用。\cmt{F6-A01}{正式勘误见EGA IV第四部分印刷第347页、上传PDF第346页：要求删除（$0_{\mathrm{IV}}$，19.3.12），理由是证明错误，$W_\lambda,U_\lambda,V_\lambda$ 不是线性流形。本书书末的相应作者勘误（见合订本书末）已选译。}

设 $A$ 为拓扑环，$B$ 为形式光滑 $A$-代数，$C$ 为拓扑 $A$-代数，$\mathfrak j$ 为其理想，满足：

$1^\circ$ $C$ 有开理想基本系 $(\mathfrak a_\lambda)$，使 $C_\lambda=C/\mathfrak a_\lambda$ 为Artin环，且典范同态 $C\to\varprojlim_\lambda C_\lambda$ 为拓扑环同构；

$2^\circ$ $\mathfrak j$ 在 $C$ 中闭且拓扑幂零；

$3^\circ$ $B$ 的每个开理想的平方仍开。

则每个连续 $A$-同态 $u:B\to C/\mathfrak j$ 分解为 $B\xrightarrow{v}C\to C/\mathfrak j$，其中 $v$ 连续。

令 $\mathfrak j_\lambda=(\mathfrak j+\mathfrak a_\lambda)/\mathfrak a_\lambda$，它是 $C_\lambda$ 的幂零理想；令 $u_\lambda$ 为 $B\xrightarrow{u}C/\mathfrak j\to C_\lambda/\mathfrak j_\lambda$ 的复合。先证明任意使 $u$ 分解为 $B\xrightarrow{w}C\to C/\mathfrak j$ 的 $A$-同态 $w$ 必连续。事实上，对每个 $\lambda$，存在 $B$ 的开理想 $\mathfrak k$，使 $u(\mathfrak k)\subset(\mathfrak j+\mathfrak a_\lambda)/\mathfrak j$。于是 $w(\mathfrak k)\subset\mathfrak j+\mathfrak a_\lambda$，故某个幂 $\mathfrak k^{2^m}$ 满足 $w(\mathfrak k^{2^m})\subset\mathfrak a_\lambda$。因 $\mathfrak k^{2^m}$ 开，连续性成立。因此只须寻找这样的 $A$-同态 $w$，不必再顾虑连续性。

原证明继续说：全体 $A$-代数同态满足
\[
 \Hom(B,C)=\varprojlim_\lambda\Hom(B,C_\lambda),
\]
而所考虑的同态集合闭于配备乘积拓扑的 $C$-模 $C_\lambda^B$ 中；由于 $C_\lambda$ 为Artin环，该乘积线性紧。对每个 $\lambda$，令 $W_\lambda$ 为使 $u_\lambda$ 分解为 $B\xrightarrow{w_\lambda}C_\lambda\to C_\lambda/\mathfrak j_\lambda$ 的 $w_\lambda\in\Hom(B,C_\lambda)$ 的集合。原证明称 $W_\lambda$ 是 $C$-模 $\Hom(B,C_\lambda)$ 中的闭线性流形，并由 $B$ 形式光滑得其非空。在乘积
\[
 E_\lambda=\prod_{\mu\leq\lambda}\Hom(B,C_\mu)
\]
中，考虑满足 $w_\mu=\varphi_{\mu\lambda}\circ w_\lambda$（$\mu\leq\lambda$）且 $w_\lambda\in W_\lambda$ 的族构成的所谓闭线性子流形 $U_\lambda$；其中 $\varphi_{\mu\lambda}:C_\lambda\to C_\mu$ 为典范同态。最后，令 $V_\lambda$ 为全乘积 $\prod_\mu\Hom(B,C_\mu)$ 中，由 $U_\lambda$ 和所有 $\mu\nleq\lambda$ 对应因子的乘积形成的所谓闭线性流形。\cmt{F6-G03}{“模”“线性流形”等用语在此是被撤销原证明的论断，不是本译稿的认定。代数同态的乘法约束非线性；报告给出平方零以外的幂零增厚例子，直接说明固定剩余同态的提升集合也未必为仿射子空间。}
原证明把问题归结为 $\bigcap_\lambda V_\lambda\ne\varnothing$，因为交中的元素按定义属于 $\varprojlim_\lambda\Hom(B,C_\lambda)$。此外，原证明使用 $\mathfrak j$ 闭以及 $C=\varprojlim C_\lambda$ 线性紧，来识别 $C/\mathfrak j\simeq\varprojlim C_\lambda/\mathfrak j_\lambda$，并把 $\varphi:C\to C/\mathfrak j$ 识别为 $\varprojlim\psi_\lambda$，其中 $\psi_\lambda:C_\lambda\to C_\lambda/\mathfrak j_\lambda$ 为典范映射，继而声称如（19.3.10）得到 $\varphi\circ w=u$。\par
\noindent\textbf{〔撤销条目译文结束。〕}

